%%%%%%%%%%%%%%%%%%%%%%%%%%%%%%%%%%%%%%%%%%%%%%%%%%%%%%%%%%%%%%%%%%%%%%%%%%%%%%%%%%%%%%%%%%%%%%%%
%% The statistical analysis: the plan as registered, then applied. Written 28 November 2026.
%%
%% Every threshold quoted here is in protocols/PROTOCOL-common.md and was committed before the
%% quantity it judges existed. The commits are named in the text on purpose: for this journal the
%% audit trail is part of the method, not an appendix to it.
%%
%% Sources: runs/task_level/{per_slot,pooled,variance_components,moderator,equivalence}.*
%%          runs/across_cohorts/* for the corpus-level layer
%%%%%%%%%%%%%%%%%%%%%%%%%%%%%%%%%%%%%%%%%%%%%%%%%%%%%%%%%%%%%%%%%%%%%%%%%%%%%%%%%%%%%%%%%%%%%%%%

\section{Statistical analysis}
\label{sec:statistics}

The numbers of \cref{sec:results} are judged here, against thresholds that were written down before
those numbers existed. We give the plan first in each subsection and then apply it, because for
this study the order matters: the plan is version-controlled and dated, and the only reason its
verdicts can be taken at face value is that it could not have been written after the fact.

\subsection{How the plan was registered}
\label{sec:registration}

Each corpus has a protocol document fixing its inclusion rule, task definition, contact rule, unit
conversion and split design, committed before any classifier was fitted on that corpus. Where the
study's own results forced a question the protocol had not anticipated, the answer is a dated
addendum committed \emph{before} the run it governs, never a rewrite: addendum~B declared the
within-corpus task contrast of \cref{sec:results-diagramo}; C the control arms under the backbone;
E the null-channel probe; F the per-descriptor baseline; G the split-noise measurement; H the
corpus-level layer; and I the task-level layer of this section, committed as \texttt{4c3eb0e} at a
point when no per-slot quantity on the confirmation splits had been computed and the script that
computes them did not exist.

Three disclosures are part of the registration rather than of the discussion, because they bound
what the verdicts below can mean, and all three were written into addendum~I before it ran. First,
the direction of the task property is \emph{not} a blind prediction: it is this study's own
corpus-level ordering, restated as a model so that it could fail. Second, the two task sets of
DiaGraMo had already been compared under addendum~B and that comparison had already failed, so
addendum~I recorded in advance that H15 --- the null --- was the outcome it expected. Third, PaHaW's
tasks are not blind, since the prior single-corpus paper reported its endpoint over a subset of
slots; PaHaW is therefore carried in the primary analysis and removed again in a declared
sensitivity analysis.

\subsection{Within a corpus}
\label{sec:stats-within}

Intervals are $10{,}000$ label-stratified participant bootstrap resamples; paired contrasts use
participant-wise swap randomisation; all repeats of a participant move together, so no resample
splits a participant across arms. These are the same estimators, with the same seeds, as the prior
single-corpus study, and the task-level layer reuses them unchanged --- which is what makes its
regression control meaningful: the participant-level margin re-derived from the task-level
prediction files reproduces the archived value in all twelve corpus-by-descriptor cells to within
$5 \times 10^{-5}$.

\subsection{Equivalence, instead of a null result}
\label{sec:stats-equivalence}

\paragraph{The plan.} A margin that lands inside an interval containing zero is not evidence of
equivalence; it is also what too few participants look like. We therefore declared a margin
$\delta = 0.03$ macro F1 and two one-sided tests at $\alpha = 0.05$, which read the $90\,\%$
interval of a slot's margin \citep{schuirmann1987,lakens2018equivalence}. Each slot receives exactly
one of four outcomes: \emph{superior} when the lower bound clears $+\delta$, \emph{inferior} when
the upper bound falls below $-\delta$, \emph{equivalent} when the whole interval lies inside
$\pm\delta$, and \emph{inconclusive} otherwise. ``Inconclusive'' is a statement about precision and
is never reported as equality.

The margin is not a convenience. It is the band this project has used since the single-corpus work
to mean ``within split noise'', and \cref{sec:results-splitnoise} measures the half-widths that
justify it. One value is used for all thirty slots, because a margin taken from each slot's own
noise would make the threshold depend on the data it judges.

\paragraph{Applied.} Of the thirty slots, two are superior, six are equivalent, twenty-two are
inconclusive and none is inferior. \textbf{H16} required at least $8$ of the $15$
shape-informative slots to be declared equivalent; $4$ are, so \textbf{H17 holds}. The claim that
the encoding is demonstrably as good as the bare trace where the shape carries the label is
\emph{not} supported at this margin --- not because the margins there are large, but because eleven
of those fifteen slots cannot resolve $0.03$: twenty-six of the thirty slots have an
MDE$_{80}$ above the margin itself. The four slots that do support it are the three character
pairs, whose standard errors are an order of magnitude smaller than any clinical slot's, and the
copy slot of DiaGraMo.

This is the first place in this study where an equivalence claim was available to be made, and the
answer is that the data do not carry it. Every ``inside the $\pm 0.03$ band'' sentence in the
earlier framing of this work was a power statement.

\subsection{Pooling thirty slots nested in six corpora}
\label{sec:stats-pooling}

\paragraph{The plan.} Slot margins are pooled by random effects fitted with restricted maximum
likelihood and reported with a Hartung--Knapp interval \citep{hartung2001}, fixed in advance because
the method-of-moments estimator is biased downward at this number of studies
\citep{viechtbauer2005}. We report $\hat\tau$, $I^2$ \citep{higgins2002} and a prediction interval
beside every pooled point, because a pooled point summarises an effect only if there is one effect.
Thirty slots measured on six groups of people are not thirty independent units, so heterogeneity is
decomposed: the same estimator with an intercept gives the total, and with corpus fixed effects ---
which absorb every between-corpus difference --- gives what remains between slots of one corpus.

\paragraph{Applied.} Over the thirty slots the pooled margin is $+0.000$, interval
$[-0.008, +0.009]$, with $\hat\tau = 0.0118$ and $I^2 = 53\,\%$. Neither sensitivity analysis moves
it: $-0.000$ without PaHaW, $+0.000$ without the corpus whose floor exceeds its best real arm.
Under the frozen backbone the thirty slots are one effect of $+0.000$ with $\hat\tau = 0$.

The decomposition is the result of this subsection, and it was measured rather than tested, as
addendum~I declared. Between slots of one corpus $\hat\tau = 0.0098$; between corpora
$\hat\tau = 0.0067$. \textbf{Task-to-task variation inside a corpus is larger than
corpus-to-corpus variation}, and both are smaller than the median sampling standard deviation of a
single slot, $0.0164$. The heterogeneity that the corpus-level layer attributed to differences
between corpora ($\hat\tau = 0.031$, $I^2 = 65\,\%$) does not survive being asked at which level it
lives.

One consequence of the change of unit must be stated and was not registered, so we flag it as an
unregistered observation: pooling slots weights a corpus by how many tasks it administered, so
DiaGraMo contributes fourteen units and Drot\'ar--Dobe\v{s} one. The slot-level pooled estimate of
$+0.000$ and the corpus-level $+0.017$ are therefore not competing estimates of one quantity; they
are two different estimands, and which one answers a question depends on whether the question is
about corpora or about tasks.

\subsection{The task property, tested within corpora}
\label{sec:stats-moderator}
%% Alias for the labels the shared generated tables reference. tables/*.tex are emitted by
%% scripts/x3_paper_tables.py for all three builds and their captions cite the earlier
%% draft's section names; aliasing here keeps one generator for three documents instead of
%% branching it per document.
\label{sec:when}

\paragraph{The plan.} The corpus-level test of this property is bounded by arithmetic: with six
corpora and two on one side, the exact permutation cannot return a $p$ below $1/15 = 0.067$, which
addendum~H stated in advance and never called significance. The task-level test is not bounded that
way. The property is entered in a meta-regression \emph{with corpus fixed effects}, so the
coefficient is a within-corpus contrast and no between-corpus difference can reach it, and it is
tested by a \textbf{restricted permutation} that holds each corpus's number of figure-imposed slots
fixed and re-labels only inside corpora. Only PaHaW, one of eight, and DiaGraMo, ten of fourteen,
vary under it, giving $\binom{8}{1}\binom{14}{10} = 8{,}008$ assignments, enumerated exactly, with
an attainable minimum $p$ of $0.00012$. The weights are held at the fitted model's, which is the
standard permutation of a meta-regression coefficient and is recorded as an implementation choice.
\textbf{H14} required both a coefficient $\ge +0.05$ and $p \le 0.05$.

\paragraph{Applied.} The coefficient is $-0.0090$ and $3{,}128$ of the $8{,}008$ assignments are at
least as extreme as the observed one, exact $p = 0.391$. \textbf{H15 holds}, in both of its
conditions and with the sign reversed from the prediction. Once the corpus is held fixed, whether
the task imposes the figure does not account for the slot-level margins. The unrestricted
permutation, which mixes the between-corpus differences back in, gives $-0.0103$ and $p = 0.158$;
it is reported for contrast and was declared as such.

The two sides of the property differ, but not in the way the hypothesis named: the fifteen
figure-imposed slots are \emph{heterogeneous} around zero ($-0.001$, $\hat\tau = 0.0212$,
$I^2 = 65\,\%$) while the fifteen shape-informative slots are \emph{homogeneous} and slightly
positive ($+0.005$, $\hat\tau = 0.0055$, $I^2 = 17\,\%$). Where the figure is imposed the encoding
sometimes does a great deal and usually nothing; where the shape carries the label it reliably does
almost nothing. That is a different statement from the one we registered, it is not a test, and
\cref{sec:discussion} treats it as the hypothesis for another study rather than as a result of this
one.

\subsection{Multiplicity}
\label{sec:stats-multiplicity}

\paragraph{The plan.} Holm correction \citep{holm1979} within families declared in the protocol.
The thirty slot contrasts are one such family, declared in addendum~I, and are corrected as one ---
the family this study actually ran, rather than one family per corpus.

\paragraph{Applied.} Five of the thirty randomisation $p$ values are below $0.05$ uncorrected.
\textbf{After Holm over the thirty, exactly one survives, and it is the smallest margin of the
five}: the \textit{m}/\textit{n} character pair, $+0.020$, adjusted $p = 0.003$ --- a contrast whose
magnitude is inside the equivalence margin. The two slots the equivalence procedure calls superior
do not survive: the static spiral of ParkinsonHW, $+0.152$, has an adjusted $p$ of $0.123$, and the
Zhu spiral, $+0.119$, of $0.549$.

The pairing is worth naming, because it is the whole argument of this section in two rows of a
table. The largest effects in the study are the least certain, and the only certain one is too small
to matter. The superiority verdicts in \cref{tab:tasklevel} are uncorrected by construction --- the
equivalence procedure reads an interval, not a corrected $p$ --- and a reader who wants a
multiplicity-controlled claim about those two spirals does not have one here.

\subsection{Power, reported with every null}
\label{sec:stats-power}

\paragraph{The plan.} For every unit we report MDE$_{80}$, the margin it could have detected at
$80\,\%$ power and two-sided $5\,\%$ given its own standard error. A null above a unit's
MDE$_{80}$ is an absence of evidence; a null below it is evidence of absence. Without the column a
reader cannot tell which one they are looking at, and this study has published both kinds.

\paragraph{Applied.} At corpus level, five of the six corpora could not have detected the pooled
margin of $+0.017$; only Character Trajectories, with $853$ recordings of one writer, could. At slot
level, twenty-six of the thirty slots could not have detected the equivalence margin of $0.03$. The
eight PaHaW slots need between $0.055$ and $0.108$; the two ParkinsonHW slots $0.142$ and $0.127$;
the two Zhu slots $0.136$ and $0.096$. The clinical corpora in this literature, at the sizes they
are released in, are not instruments that can resolve effects of the size this recipe produces.

\subsection{Screening and confirmation as a two-stage design}
\label{sec:stats-twostage}

\paragraph{The plan.} The two master seeds form a two-stage design whose interim rule was fixed
before either stage ran: a representation reaches confirmation only if it leads at screening, and no
margin is reported as a gain unless it reappears on the confirmation seed inside the declared band.
We claim \emph{no} propagation of $\alpha$ between the stages --- they are separate families and the
confirmation intervals are not widened for having passed a screen.

\paragraph{Applied.} What the screen can buy is measured instead of argued: picking the arm on the
split it is scored on is worth up to $+0.028$ on one corpus and nothing on five. Every number in
\cref{sec:results} is from the confirmation stage, and the task-level layer draws no new split at
all --- it re-slices the confirmation predictions, so it inherits that stage's separation rather
than creating a third one.

\subsection{What the registrations predicted, and what happened}
\label{sec:stats-ledger}

Nine hypotheses were registered with a direction and a threshold, across addenda~B, C, E, F, H and
I. \textbf{Three held.} H6, the null-channel probe: the gain on Zhu is the channel's content and not
the mere presence of a varying channel. H10, heterogeneity between corpora under LPQ. H12, the
corpus-level moderator --- at a $p$ equal to its attainable minimum, which is why it is reported as
an extremity and not as significance.

\textbf{Six failed.} H1 and H2, the within-corpus task contrast on DiaGraMo, with the sign reversed.
H4, the spatial controls under the frozen backbone, where the test dissolved its own target because
the backbone had no margin left for a control to destroy. H8, that at least half of a reported
margin survives a descriptor-matched baseline --- refused on all three corpora that had a margin to
decompose. H14, the task property within corpora. H16, equivalence on the shape-informative slots.

We report the ledger rather than only its favourable rows because the ledger is the result. A study
of this shape earns its conclusions from the registrations that failed, and this one has more of
those than of the other kind.
